\section*{結言}
今回の観測窓は、大型モデルの新発表が相次いだ直後に当たり、Mistral Large 4とGemini 4 Carbonへの反応や比較が中心となった。オープンウェイトと限定提供という公開形態の違いが、自己ホストの可能性、利用条件、サイバー防御への適用という実務的な論点を浮かび上がらせている。性能に関する投稿上の主張と、再現可能な評価結果は分けて読む必要がある。

共有robots.txt信号の議論は、モデルの能力向上と同時に学習データの利用制御が制度・実装上の課題となっていることを示した。Anthropicの自動レッドチーミングとPlanAssistの自治体展開は、安全性検証と現場の業務改善という異なる方向から、エージェント技術の実用化を映している。

モデルの更新速度を追う投稿やOpenCodeの運用談義には、利用者が多様な選択肢を比較し続ける状況が表れた。次に重要になるのは、各社の実装時期とアクセス条件、現場での効果測定、独立した性能再現の蓄積である。本号はGrokの観測結果を記事化したもので、追加の一次資料検証は行っていない。
